\section{How the Church learned to keep Advent}

\subsection{A date in Rome, a congregation in Spain}

The first name in a Roman list of martyrs' anniversaries is Christ's. The list
preserved in the \emph{Chronography of 354} records his birth at Bethlehem on
the eighth day before the Kalends of January. Count as a Roman did, including
both ends, and the date is December 25. Farther down the page are the places
where Romans kept the anniversaries of Agnes, Lawrence and other martyrs.
Here, in a congregation's calendar, is evidence that Christians at Rome already
knew when to keep the Lord's birthday. To learn how they prepared for it one
has to look elsewhere.\footnote{\emph{Depositio Martirum}, opening entry,
ed. Theodor Mommsen, \emph{Chronica minora} I (Berlin, 1892),
\href{https://www.dmgh.de/mgh_auct_ant_9/index.htm\#page/71/mode/1up}{p.~71}.}

Farther west, bishops gathered at Zaragoza in 380 told the faithful where
to spend the twenty-one days beginning December 17. They were to come to
church together. The bishops expressly mention houses, country estates and
mountains as places where people might withdraw. The period included Christmas
and continued through January 6, the Epiphany. A reader seeking the first
four-Sunday Advent in this canon will find a different concern. The bishops
wanted Christians who might otherwise make a private observance to worship
with their neighbors, and their instruction crossed the boundary that now
separates Advent from Christmas Time.\footnote{First Council of Zaragoza,
canon~4, in G.~D. Mansi, \emph{Sacrorum conciliorum nova et amplissima
collectio}, III (Florence, 1759), cols.~633--635.}

In fifth-century Rome, St Leo spoke to people who had finished bringing in
the crops. A full storehouse could teach its owner to thank God, to eat less
and to make more food available to people who had none. Leo preached on
Christ's warning in Luke 21 that anxiety and excess could weigh down the
heart. In December it made sense to ask both whether the heart was awake
and what had become of the harvest. He sets this fast alongside the fasts of
spring, summer and autumn. His sermons tell us what he wanted his congregation
to do in winter. The Roman Sunday cycle can be traced in later books.\footnote{Leo
the Great, Sermons 16.2 and 19.1--3, trans. C.~L. Feltoe,
\emph{Nicene and Post-Nicene Fathers}, second series, XII (1895),
\href{https://www.ccel.org/ccel/schaff/npnf212.html}{Sermons on the December fast}.}

\subsection{Counting back from Christmas}

In Tours, the feast that marked the change of diet was the burial feast of the
city's own bishop, St Martin, on November 11. Gregory of Tours, writing in the
sixth century, records the arrangements of his fifth-century predecessor
Perpetuus. From October until Martin's feast there had been two fasts a week.
From Martin's feast to the Nativity there were three. At the Council of Tours
in 567 the bishops prescribed a different rhythm for monks, who were to fast
daily from the beginning of December to Christmas. Another council, at Mâcon
later in the sixth century, named the three days of the week between St
Martin's and Christmas on which the faithful were to fast. Monday, Wednesday
and Friday were to have Mass in the Lenten order. The season sometimes called
St Martin's Lent has a recognizable face here, down to the days on which a
household changed its meals. Its extent and demands varied even within Gaul.
\footnote{Gregory of Tours, \emph{Histories} X.31, in O.~M. Dalton,
\emph{The History of the Franks}, II (Oxford, 1927), pp.~471--472.
Second Council of Tours, c.~18 (17), and First Council of Mâcon, c.~9,
in F.~Maassen, ed., \emph{Concilia aevi Merovingici} (Hanover, 1893),
\href{https://www.dmgh.de/mgh_conc_1/index.htm\#page/126/mode/1up}{p.~126}
and \href{https://www.dmgh.de/mgh_conc_1/index.htm\#page/157/mode/1up}{p.~157}.}

A sacramentary gives us another way into the season. It is the book of
prayers the priest said at Mass. The eighth-century manuscript known as the
Old Gelasian Sacramentary has five groups of Advent Mass prayers. The later
headings introduce another Mass without assigning a number to a Sunday.
In the second group the priest asks God to rouse the people's hearts to
prepare the Son's way. He goes on to ask for grace that will go before them
and follow them. A third prayer asks God to visit them daily and cleanse their
consciences, so that the Son will find a home ready for him. The one for whom
the house is being prepared is already at work there. Five sets of prayers are
secure. The association of each set with a particular Sunday has to be
inferred by comparison with other books.\footnote{H.~A. Wilson, ed.,
\emph{The Gelasian Sacramentary} (Oxford, 1894), book II, nos.~80--84,
especially \href{https://archive.org/details/LiberSacramentorumRomanaeEcclesiae}{pp.~214--215 and p.~214 n.~1}.}

When Charlemagne asked Pope Hadrian I for a Roman sacramentary, another
arrangement travelled north. The Gregorian book sent from Rome in the late
eighth century is known through later copies. Its first Advent Sunday asks
the Lord to stir up his power and come. The third is assigned to St Peter's,
the Roman station or church for the day's Mass. After the December Ember
Days, seasonal days of prayer and fasting, a heading marks the Sunday as
vacant. There is nevertheless a prayer beneath it asking the Lord to come
quickly and overcome the delay our sins have caused. The ninth-century
writer Amalarius of Metz had good reason to puzzle over the number of weeks.
He could find readings extending back five Sundays, while the Missal and
antiphoner he discussed called for four weeks of preparation. His explanation
of the four reaches for the Law, the Prophets, the Psalms and the opening of
the Gospel. He was reading the one promise of Christ across books whose
calendars still differed.\footnote{H.~A. Wilson, ed., \emph{The Gregorian
Sacramentary under Charles the Great} (London, 1915), introduction p.~xxi,
\href{https://commons.ptsem.edu/id/gregoriansacrame0049cath}{nos.~150--157,
pp.~113--116}. Amalarius, \emph{Liber officialis} III.40, \emph{Patrologia
Latina} 105, cols.~1158--1159.}

\subsection{What the choir kept and the printers handed on}

In the antiphoner of the Benedictine abbey of St Gall, written around the
year 1000, a singer could follow a series of great invocations for Vespers
in the week before Christmas. Wisdom, Lord, Root of Jesse, Key of David,
Dayspring, King of the nations, Emmanuel. Each title gives the choir a fresh
reason to ask Christ to come. The sequence fills one page and turns onto
the next. St Gall then adds more chants, including an imagined conversation
in which the daughters of Jerusalem question the Virgin. Twelve pieces and
a heading for the week tell us what this choir could sing. The seven dates
familiar in the Roman Office today come from that Office's own order.
\footnote{St Gall, Stiftsbibliothek, Cod.~Sang.~390, pp.~40--41, reproduced
in \emph{Paléographie musicale}, second series, I (Solesmes, 1900).
\href{https://cantusdatabase.org/chant/456197/}{Manuscript chant index}.
United States Conference of Catholic Bishops, \href{https://www.usccb.org/prayers/o-antiphons-advent}{The O Antiphons of Advent}.}

By the late thirteenth century William Durandus, bishop of Mende, could
take for granted a musical experience familiar to Catholics today. The choir
sings Alleluia before the Advent Sunday Gospel after passing over the Gloria
at the beginning of Mass. Durandus heard joy in the coming already accomplished
and sorrow for sin in the thought of judgment. He heard something else in the
silence, too. When the angels' hymn returned on the night of the Nativity,
a congregation that had waited for it could sing it with longing. His book
also attributes the four weeks to St Peter. That attribution belongs to a
medieval reading of the season. The early books we possess preserve a more
varied history.\footnote{William Durandus, \emph{Rationale divinorum
officiorum}, VI.2, \href{https://archive.org/details/rationalediuinor02dura}{Lyon printing of 1612},
II, ff.~255v--256v.}

Open the Roman Missal that Pius V promulgated after Trent in the Venice
printing of 1570. The First Sunday has its station at St Mary Major, the
entrance from Psalm 24 (25 in modern numbering), the prayer asking the Lord
to stir up his power, and Luke's Gospel of the Son of Man coming in a cloud.
The 1962 Roman Missal retains this combination. Both books carry the Ember
Masses after the third Sunday. They lead from Gabriel's greeting of Mary on
Wednesday to her visit to Elizabeth on Friday. Comparing the two Missals
also reveals change within the older Roman order. The printed rubric on
the first Sunday in 1570 links the dismissal to the omission of the Gloria.
In the ordinary circumstances of 1962 the usual dismissal was already
restored. Rome's books coexisted with other Western uses. In an Ambrosian
Missal printed in Milan in 1712, a reader can turn to a \emph{sixth} Advent
Sunday and find a Mass of the Incarnation along with the Mass of Advent.
\footnote{\emph{Missale Romanum}, \href{https://archive.org/details/messale-tridentino-di-san-pio-v-1570}{Venice, 1570},
Advent I--IV, especially ff.~1r--2r and 8v--9r. \emph{Missale Romanum},
\href{https://musicasacra.com/pdf/missale62.pdf}{Vatican typical edition, 1962},
pp.~1--16. General Rubrics (1960), no.~507. \emph{Missale Ambrosianum},
\href{https://archive.org/details/10869665bsb}{Milan, 1712}, pp.~12--15.}

\subsection{What changed in the parish church}

In 1963 the Second Vatican Council asked for a larger share of Scripture
to be heard at Mass over a cycle of years. Its reason reaches beyond the
number of readings. In the course of a year the Church unfolds the mysteries
of Christ so that the faithful can draw on the grace of what they celebrate.
Paul VI's constitution \emph{Missale Romanum} of 1969 explained the changes
to the Order of Mass and the new three-year cycle of Sunday readings. The
first typical edition of the revised Latin Missal appeared in 1970. The
English Missal used in this handbook belongs to a later stage, the U.S.
translation of the third typical edition, introduced in full on the First
Sunday of Advent, November 27, 2011. The Missal gives the priest the prayers.
The Lectionary gives the readings proclaimed at the ambo.
\footnote{Second Vatican Council, \href{https://www.vatican.va/archive/hist_councils/ii_vatican_council/documents/vat-ii_const_19631204_sacrosanctum-concilium_en.html}{\emph{Sacrosanctum Concilium}},
nos.~51, 102. Paul VI, \href{https://www.vatican.va/content/paul-vi/en/apost_constitutions/documents/hf_p-vi_apc_19690403_missale-romanum.html}{\emph{Missale Romanum}},
April 3, 1969. USCCB, \href{https://www.usccb.org/prayer-and-worship/the-mass/general-instruction-of-the-roman-missal/roman-missal}{Roman Missal}
and \href{https://www.usccb.org/resources/divine-worship-newsletter-aug-sept-2010}{\emph{Divine Worship Newsletter}, August--September 2010}.}

Listen at the two ends of Advent. The prayer at the opening of the old
First Sunday asked the Lord to come to his endangered people. The collect
in the revised Missal asks the Father to give us the will to go out to meet
his Son. It pictures those who bring deeds of justice finding a place at
Christ's right hand, an image from Matthew 25. On the Fourth Sunday the
revised Missal takes up the prayer many Catholics know from the Angelus.
We have heard the angel's tidings of the Incarnation. The priest asks that
grace will carry us through the Lord's Passion and Cross to his risen glory.
On the Sunday nearest the manger, the prayer already names Easter.
\footnote{\emph{Roman Missal}, Third Edition, Proper of Time, Collects of the
First and Fourth Sundays of Advent, in the USCCB's
\href{https://www.usccb.org/prayer-and-worship/liturgical-year-and-calendar/advent/commentary-on-advent-proper-prayers}{Commentary on the Proper Prayers of Advent}.}

The Sunday Gospels retain the First Sunday's watch for the Lord's return.
In Year C one can still hear the image of the cloud in Luke 21. On the next
two Sundays John the Baptist bears witness. The Fourth Sunday brings us to
Joseph hearing the angel in Year A, to Gabriel and Mary in B, and to Mary
and Elizabeth in C. Some of the scenes once heard at the December Ember
Masses now reach the congregation on a Sunday. Weekday Mass, too, has its
own course. Earlier in Advent Isaiah supplies a series of readings. From the
Thursday of the second week the Gospel puts John the Baptist in view. The
December 17--24 series draws on Matthew 1 and Luke 1 and pairs them with
passages from several Old Testament books. The revised calendar leaves
conferences of bishops to arrange Ember and Rogation observances according
to local need. In the United States the diocesan bishop sets these days of
prayer. An old missal's three Ember Masses no longer tell a reader
which Gospel will be proclaimed in the parish that week.
\footnote{\emph{Lectionary for Mass}, General Introduction,
\href{https://www.liturgyoffice.org.uk/Resources/GIRM/Documents/Lectionary.pdf}{nos.~93--94}.
U.S. Lectionary, nos.~1--12, 175--191, 193--200. \emph{Universal Norms on the
Liturgical Year and the Calendar}, nos.~39--42, 45--47. USCCB,
\href{https://www.usccb.org/committees/divine-worship/special-observances}{Days of Prayer and Special Observances}, citing GIRM no.~373.}

Paul VI promulgated the revised Liturgy of the Hours in 1970. Morning and
Evening Prayer stand at the two hinges of the day. The psalms are arranged
over four weeks. The Office of Readings joins an extended passage of Scripture
to a reading from the Church's tradition and can be prayed during the day
or kept as a night watch. At Vespers from December 17 through 23, the seven
great invocations still accompany Mary's Magnificat. A Christian who sings
one of them today shares an address to Christ with those medieval singers at
St Gall. The season begins at First Vespers of the Sunday on or nearest
November 30 and ends before First Vespers of the Nativity. The prayers of
two successive evenings can take a household all the way across that
threshold. On December 23 the Church calls for Emmanuel, and on Christmas
Eve she begins to celebrate his birth.
\footnote{Paul VI, \href{https://www.vatican.va/content/paul-vi/la/apost_constitutions/documents/hf_p-vi_apc_19701101_laudis-canticum.html}{\emph{Laudis canticum}},
November 1, 1970, especially nos.~2--5; \emph{Liturgy of the Hours},
vol.~I, Proper of Seasons, December 17--23, Evening Prayer.
\href{https://www.usccb.org/prayer-worship/liturgical-year/advent}{Universal Norms},
nos.~39--42.}
